% chapters/ch34_r1fix_patch.tex
% Chapter 34 (Part VI): R1-FIX Patch — Revised Chapters 26-35 of v3.0
% Source: NRMO_v5_updated.pdf, Part VI Ch33 (lines 10595-10860 of v5_layout.txt)

\chapter{R1-FIX Patch --- Revised Chapters 26--35
(NRMO v3.0 Specification)}
\label{ch:r1fix-patch}

\nrmobox{Document status}{
\textbf{Title}: NRMO v3.0 specification | Revised R1-FIX Chapters
26--34\\[0.3em]
\textbf{Scope}: 5 chapters | Parallel + Design\\[0.3em]
\textbf{Status}: REVISION DRAFT --- Pending Constitutional Review\\[0.3em]
\textbf{Affects}: NRMO v3.0 (Chapters 26--34); does \textbf{not}
modify v2.0 design (parallel + CONSTITUTIONAL FROZEN).
}

% =====================================================================
\section{Conflict Identifiers Resolved}
\label{sec:r1fix-conflicts}

\begin{table}[ht]
\centering
\small
\begin{tabularx}{\linewidth}{lXl}
\toprule
\textbf{ID} & \textbf{Conflict description} & \textbf{Chapter} \\
\midrule
C-01 & Parallel execution vs.~Type ZERO suppression. & Ch.~27 \\
C-04 & APSOC definition state DEADLOCK. & Ch.~30 \\
D-01 & Scoreboard parameter weighting. & Ch.~31 \\
D-02 & $\varepsilon_{\max}(B_t)$ definition. & Ch.~26 \\
D-03 & Challenge Window update condition. & Ch.~31 \\
\bottomrule
\end{tabularx}
\caption[R1-FIX conflicts]{Conflicts identified for R1-FIX revision
of NRMO v3.0 Chapters 26--35.}
\label{tab:r1fix-conflicts}
\end{table}

% =====================================================================
\section{Chapter 26 Invariants (Revised)}
\label{sec:r1fix-invariants}

The following invariants are FROZEN under the R1-FIX revision:

\begin{itemize}
\item \textbf{NRMO CORE / Ruin}: Ruin is the absorbing failure
state; $\varepsilon$ bounds action probability of entering Ruin.
\item \textbf{Shutdown Atomicity}: Shutdown is atomic; partial /
interrupted execution prohibited.
\item \textbf{Type ZERO suppression}: Type ZERO exercises
suppression in Phased Parallel Phase~3 (Section~27.3); Phase~1
evaluation is read-only.
\item \textbf{DAG structure}: Definition \& Assertion Governance;
directed acyclic consistency.
\item \textbf{Non-Personification}: Modules are organs, not agents.
\item \textbf{$A_{\mathrm{allowed}}$ constraint}: Permitted action
set externally supplied; NRMO does not generate or modify.
\item \textbf{APSOC Pre-Commit Gate}: Execution requires
reversibility / loss / condition gates passed.
\item \textbf{$\varepsilon$-Supply}: $\varepsilon$ supplied
externally; $\varepsilon_{\max}(B_t)$ constraint
(Section~26.1 of v3.0).
\end{itemize}

% =====================================================================
\section{$\varepsilon_{\max}(B_t)$ Definition}
\label{sec:r1fix-epsilon-max}

\nrmobox{Specification warning}{
The specification ``$\varepsilon_{\max}(\cdot)$'' must be carefully
distinguished from the institutional $\varepsilon$
(Chapter~\ref{ch:epsilon-supply}). $\varepsilon_{\max}(B_t)$ is a
buffer-modulated dynamic ceiling, \emph{never} above the
institutional $\varepsilon$.
}

\subsection{$B_t$ Buffer State}

$B_t$ is the buffer state, normalised to $[0, 1]$:

\begin{equation}
B_t = \frac{B_{\mathrm{capital}}}{B_{\mathrm{max}}}
\;\in\; [0, 1].
\end{equation}

\noindent
Or alternatively:
\begin{equation}
B_t = \frac{\mathrm{remaining}}{\mathrm{initial}}
\;\in\; [0, 1].
\end{equation}

\noindent
For multi-resource composition:
\begin{equation}
B_t = \sum_i w_i \cdot \frac{b_i}{b_{i, \max}},
\quad B_t \in [0, 1].
\end{equation}

\paragraph{Constraint.}
$B_t \in [0, 1]$. When $B_t = 0$: SHUTDOWN-equivalent.

\subsection{$\varepsilon_{\max}(B_t)$ Selection Conditions}

The $\varepsilon_{\max}$ function must satisfy:

\begin{enumerate}
\item \textbf{Monotonicity}: $B_{t_1} \le B_{t_2} \Rightarrow
\varepsilon_{\max}(B_{t_1}) \le \varepsilon_{\max}(B_{t_2})$.
\item \textbf{Floor at zero buffer}:
$\varepsilon_{\max}(0) = \varepsilon_{\mathrm{floor}} > 0$ (zero
buffer state).
\item \textbf{Ceiling at full buffer}:
$\varepsilon_{\max}(1) = \varepsilon_{\mathrm{global}}$
(institutional ceiling).
\item \textbf{Smoothness}: $\varepsilon_{\max}(\cdot) \in C^1$
(continuously differentiable).
\end{enumerate}

\subsection{Canonical $\varepsilon_{\max}$ Function}

\begin{equation}
\boxed{
\varepsilon_{\max}(B_t) =
\varepsilon_{\mathrm{floor}}
+ (\varepsilon_{\mathrm{global}} - \varepsilon_{\mathrm{floor}})
\cdot (1 - e^{-\kappa B_t}).
}
\label{eq:r1fix-epsilon-max-canonical}
\end{equation}

\paragraph{Parameter ranges.}

\begin{itemize}
\item $\varepsilon_{\mathrm{floor}}$: floor at $B_t = 0$,
typically $\varepsilon_{\mathrm{global}} / 10$.
\item $\varepsilon_{\mathrm{global}}$: institutional ceiling
($B_t = 1$).
\item $\kappa$: rate parameter; canonical $\kappa = 3.0$.
\end{itemize}

\subsection{$\varepsilon_t$ Update Protocol}

\begin{enumerate}
\item Each decision-making step updates $B_t$.
\item Compute $\varepsilon_t = \min(\varepsilon_{\max}(B_t),
\varepsilon_{\mathrm{global}})$.
\item NRMO CORE consumes $\varepsilon_t$ as the active veto
threshold.
\item All $\varepsilon_t$ changes are recorded append-only with
audit trail.
\item Edge case: $B_t = 0 \Rightarrow \varepsilon_t =
\varepsilon_{\mathrm{floor}}$, SHUTDOWN evaluation triggered.
\end{enumerate}

% =====================================================================
\section{Chapter 27 Execution Architecture (Revised)}
\label{sec:r1fix-execution-architecture}

\subsection{Design Principle: Phased Parallel Execution}

\nrmobox{Specification clarification}{
The original v3.0 specification was misread as ``parallel execution
prohibited; Type ZERO suppression overrides parallelism.'' The
correct reading: \emph{parallel execution is permitted in Phase~1
(read-only)}; Type ZERO suppression operates in Phase~3
(serial / suppression).
}

\paragraph{Three phases.}

\begin{description}
\item[Phase 1 (Fork --- parallel evaluation)]
All modules evaluate state in parallel. Read-only; no state
changes permitted. Type ZERO Phase~1: read-only evaluation.
\item[Phase 2 (Barrier --- await completion)]
Wait for all modules to complete or hit $T_{\max}$ cutoff. Late
scores set to 0.
\item[Phase 3 (Join --- serial / suppression)]
Priority Arbiter $\to$ Tri-Option Generator $\to$ APSOC $\to$
Shutdown $\to$ NRMO CORE serial. Type ZERO suppression operates
here.
\end{description}

\paragraph{Resolution.}
\begin{itemize}
\item Suppression separated from Phase~1 (evaluation) and Phase~3
(suppression).
\item Parallel evaluation and Type ZERO suppression are now
non-conflicting in time.
\end{itemize}

\subsection{Architecture Diagram}

\begin{lstlisting}[basicstyle=\ttfamily\scriptsize,frame=single]
Input (state / decision-making)
        |
        v
PHASE 1: FORK (parallel evaluation)
  [OODA] [HST-N] [Fast Strike]
  [Type ZERO*] [PP] [Heavy Strike]
  [DAG] [SOP] [Stealth] [Tuatha]
  *Type ZERO: Phase 1 = evaluation only;
   suppression in Phase 3.
        |
        v (completion or T_max)
PHASE 2: BARRIER (await)
        |
        v
PHASE 3: JOIN (serial)
  Priority Arbiter
  -> Tri-Option Generator (A/B/C)
  -> APSOC -> Shutdown (Atomic)
  -> NRMO CORE -> Action Router
\end{lstlisting}

\subsection{Steps (Numbered)}

\begin{enumerate}
\item Input: state / decision-making.
\item Phase 1 Fork: parallel module evaluation.
\item Phase 2 Barrier: $T_{\max}$ cutoff completion (overdue
modules score 0).
\item Phase 3 Type ZERO suppression: Priority Arbiter score
aggregation + execution.
\begin{itemize}
\item Tri-Option Generator: A/B/C with APSOC Pre-Commit gates.
\end{itemize}
\item Shutdown (Atomic): conditional shutdown evaluation.
\item NRMO CORE: Ruin $\le \varepsilon_t$ check.
\item Action Router: approved action execution.
\end{enumerate}

\subsection{$T_{\max}$ Cutoff}

$T_{\max}$ is the parallel-evaluation deadline. Modules whose
evaluation has not completed by $T_{\max}$ have their scores
discarded (set to 0). This prevents slow modules from blocking the
pipeline indefinitely.

% =====================================================================
\section{Boundary Contract (APSOC Gates)}
\label{sec:r1fix-boundary-contract}

Each option produced by Tri-Option Generator must pass three APSOC
gates:

\begin{itemize}
\item \textbf{Max Loss Cap}: maximum loss bounded.
\item \textbf{Reversibility Class}: reversible / partially
reversible / irreversible classification.
\item \textbf{Exit Trigger condition}: explicit state-defined exit
condition.
\end{itemize}

\subsection{A/B/C Option Allocation Constraint}

Default option-portfolio allocation:

\begin{table}[ht]
\centering
\small
\begin{tabularx}{\linewidth}{lXX}
\toprule
\textbf{Option} & \textbf{Risk class} & \textbf{Default
allocation} \\
\midrule
A (defensive) & Reversible, low-impact & 20\% (or 80\% / 30\%) \\
B (balanced) & Reversible, medium-impact & 0\% (or 80\% / 50\%) \\
C (offensive) & Boundary-approaching & 0\% (or 25\% / 20\%) \\
\bottomrule
\end{tabularx}
\caption[Boundary contract A/B/C]{Default A/B/C allocation. Constraint:
$A + B + C = 100\%$.}
\label{tab:r1fix-abc-allocation}
\end{table}

\paragraph{Constraint.}
$A + B + C = 100\%$.

% =====================================================================
\section{Scoreboard (Revised)}
\label{sec:r1fix-scoreboard}

\nrmobox{Specification clarification}{
The original v3.0 specification's scoreboard formula included
unclear parameter normalisation. R1-FIX reformulates with explicit
$[0, 1]$ normalisation per dimension.
}

\subsection{Normalised Performance Vector}

\begin{lstlisting}[basicstyle=\ttfamily\small,frame=single]
xK_DeltaL = clip(1 - DeltaLatency / DeltaLatency_ref, 0, 1)
              # Latency: 1 = best
xK_Cs     = C_success_rate
              # in [0,1]
xK_MD     = clip(1 - MaxDrawdown / MaxDrawdown_ref, 0, 1)
              # Drawdown loss: 1 = best
xK_AV     = clip(1 - (Abort + Veto) / AV_ref, 0, 1)
              # Abort+Veto: 1 = least

# Reference values DeltaLatency_ref, MaxDrawdown_ref, AV_ref
# come from A_allowed.
\end{lstlisting}

\subsection{Scoreboard Equation}

\begin{equation}
\mathrm{Score} =
w_1 \cdot \mathrm{xK}_{\Delta L}
+ w_2 \cdot \mathrm{xK}_{\mathrm{Cs}}
- w_3 \cdot (1 - \mathrm{xK}_{\mathrm{MD}})
- w_4 \cdot (1 - \mathrm{xK}_{\mathrm{AV}}).
\label{eq:r1fix-scoreboard}
\end{equation}

\noindent
Weights $w_1, w_2, w_3, w_4 \ge 0$ are tuned per Mode (S, D, M, H, T)
within $A_{\mathrm{allowed}}$.

% =====================================================================
\section{Challenge Window (Revised: Update Conditions)}
\label{sec:r1fix-challenge-window}

\nrmobox{Specification clarification}{
The original v3.0 specification stated ``A increases, C suppression''
without explicit update equations. R1-FIX provides definitive
update equations with constraint validation.
}

\subsection{Trigger Condition}

\begin{equation}
\mathrm{deterioration\_flag} = (\mathrm{Score}_t <
\mathrm{Score}_{\mathrm{baseline}}).
\end{equation}

When $\mathrm{Score}_t \ge \mathrm{Score}_{\mathrm{baseline}}$ for
consecutive periods, Challenge Window enters recovery mode.

\subsection{Update Equations}

For coefficients $\alpha, \beta > 0$:

\begin{equation}
A_{t+1} = A_t + \alpha \cdot \mathrm{deterioration\_flag} \cdot
(A_{\max} - A_t),
\label{eq:r1fix-At-update}
\end{equation}

\begin{equation}
C_{t+1} = C_t - \beta \cdot \mathrm{deterioration\_flag} \cdot
C_t,
\label{eq:r1fix-Ct-update}
\end{equation}

\begin{equation}
B_{t+1} = 1.0 - A_{t+1} - C_{t+1}.
\label{eq:r1fix-Bt-update}
\end{equation}

\subsection{Default Parameters}

\begin{itemize}
\item $\alpha = 0.20$ (A increase rate, $0 < \alpha < 1$).
\item $\beta = 0.30$ (C decrease rate, $0 < \beta < 1$).
\item $A_{\max} = 0.80$ (A ceiling).
\item $\mathrm{Score}_{\mathrm{baseline}}$: rolling N-step mean.
\end{itemize}

\subsection{Constraint Validation}

\begin{lstlisting}[basicstyle=\ttfamily\small,frame=single]
assert A_{t+1} >= 0.20            # A floor
assert C_{t+1} >= 0.0             # C floor
assert 0.0 <= B_{t+1} <= 0.80     # B range
assert abs(A_{t+1} + B_{t+1}
           + C_{t+1} - 1.0) < 1e-9
                                  # 100% allocation
\end{lstlisting}

\subsection{Recovery Condition}

\begin{equation}
\mathrm{Score}_t \ge \mathrm{Score}_{\mathrm{baseline}}
\quad \text{for } N_{\mathrm{recover}} \text{ consecutive periods}.
\end{equation}

\begin{lstlisting}[basicstyle=\ttfamily\small,frame=single]
if consecutive_recovery_count >= N_recover:
    (A, B, C) = (A_default, B_default, C_default)
              # Restore defaults
\end{lstlisting}

\subsection{Replacement and Cooling}

On shutdown (any cause), suppression cools:
$\alpha = \alpha_{\min}$, $\beta = \beta_{\max}$ for
$N_{\mathrm{cool}}$ periods.

% =====================================================================
\section{Minimal Handover Output Format}
\label{sec:r1fix-handover}

The Tri-Option Generator output for handover:

\begin{lstlisting}[basicstyle=\ttfamily\scriptsize,frame=single]
Selection {A | B | C} with Score and
Boundary Contract structure
{Max Loss, Reversibility, Exit state}.

State (A_t, B_t, C_t) and Challenge Window
consecutive_recovery_count are tracked.

Chapter 32 Output (Revised):
- Mode Weights:    {S, D, M, H, T}
- Normalized Inputs:
    x_DeltaL: [0,1]
    x_Cs:     [0,1]
    x_MD:     [0,1]
    x_AV:     [0,1]
- Options:
    A: { Ruin: __, Rev: __, MaxLoss: __,
         Expect: __, APSOC_gates: [1] }
    B: { Ruin: __, Rev: __, MaxLoss: __,
         Expect: __, APSOC_gates: [1, 2] }
    C: { Ruin: __, Rev: __, MaxLoss: __,
         Expect: __, APSOC_gates: [1, 2, 3] }
- Chosen:
    [ A | B | C | SAFE_EXIT | DEADLOCK_HOLD ]
- Reason:        [1]
- Caps:          [...]
- Score:         [...]
- Allocation:    { A: __, B: __, C: __ }
- LogID:         [audit ID]
- Emergency Bypass:
    [true | false]
    (when true: log audit + 3-step
     SHUTDOWN evaluation)
\end{lstlisting}

% =====================================================================
\section{Chapter 33 Prohibitions}
\label{sec:r1fix-prohibitions}

The R1-FIX revision lists 10 prohibited operations:

\begin{enumerate}
\item \textbf{CORE}: NRMO CORE veto-only; no override of veto.
\item \textbf{Shutdown}: Shutdown is conditional and atomic; cannot
be interrupted.
\item \textbf{Phase 1}: state changes prohibited in parallel
evaluation.
\item \textbf{Phase 3}: parallel execution prohibited; serial
sequence enforced.
\item \textbf{Invariants}: removal or weakening of any invariant
prohibited.
\item \textbf{Reason}: actions without recorded reasons prohibited.
\item \textbf{Irreversibility}: irreversible actions without explicit
exit clause prohibited.
\item \textbf{AI judgement}: AI-only decision-making prohibited.
\item \textbf{$\varepsilon_{\max}$}: changing $\varepsilon_{\max}$
function ($\varepsilon_{\mathrm{floor}}$, $\varepsilon_{\mathrm{global}}$,
$\kappa$) without governance approval prohibited.
\item \textbf{Scoreboard}: changing $w_1$--$w_4$ weights without
mode transition approval prohibited.
\end{enumerate}

% =====================================================================
\section{Chapter 34 Conditions}
\label{sec:r1fix-conditions-summary}

The R1-FIX revision lists 9 condition checks:

\begin{enumerate}
\item \textbf{Ruin condition}: $\tau_F = \infty$ (no ruin
observation).
\item \textbf{Decision-making total time}: Phase 1+3
$\le T_{\max\_\mathrm{total}}$ ($\Delta$ Latency).
\item \textbf{Firepower operation}: Stealth-Acc unit status check.
\item \textbf{Firepower state}: Phase 1+3 active modules verified.
\item \textbf{Score baseline}:
$\mathrm{Score} \ge \mathrm{Score}_{\mathrm{baseline}}$.
\item \textbf{Reason}: Evidence Capsule audit trail present.
\item \textbf{Veto / Abort}: counts recorded; AV computed.
\item \textbf{DEADLOCK}: SAFE\_EXIT delivered, DEADLOCK\_LOG
recorded.
\item \textbf{$\varepsilon_{\max}$}: decision-making
$\varepsilon_t \le \varepsilon_{\max}(B_t)$.
\end{enumerate}

% =====================================================================
\section{Chapter 35 Definition: Closure}
\label{sec:r1fix-closure}

\nrmobox{Closure declaration}{
\centering
\textit{NRMO ver3.0 = Ruin-Firepower-OS (operation).}\\[0.3em]
Revised five points:\\
parallel execution + Type ZERO $\to$ suppression separation;\\
Phased Parallel Execution;\\
APSOC + DEADLOCK;\\
Scoreboard $\to$ $\varepsilon_{\max}(B_t)$ $\to$ definition;\\
Challenge Window $\to$ update conditions.
}

\paragraph{Document closure.}
End of NRMO ver3.0 Patch R1-FIX | Full-Active Combat-Constitution OS
| R1-FIX | 2026-02-18.

% =====================================================================
\section*{Connections to Other Parts}

\begin{itemize}
\item Equation~\ref{eq:r1fix-epsilon-max-canonical} ($\varepsilon_{\max}(B_t)$)
is the buffer-modulated dynamic ceiling, used in Algorithm~1
(Section~\ref{sec:tf-v4-algorithm}) Step~1.
\item Phased Parallel Execution architecture aligns with the
Parallel OODA specification (Chapter~\ref{ch:parallel-ooda}).
\item Boundary Contract (Section~\ref{sec:r1fix-boundary-contract})
implements APSOC's Pre-Commit gate
(Section~\ref{sec:apcso-formal}).
\item Scoreboard (Section~\ref{sec:r1fix-scoreboard}) implements the
multi-axis score aggregation; the engineering implementation
(Part~X) uses an analogous scoring loop.
\end{itemize}
